The apex and the smooth point. Panels (a) and (b) draw the symmetric slice $S=\left(\begin{smallmatrix}x+y&z\\ z&x-y\end{smallmatrix}\right)$ of $2\times2$ matrices, in which rank $\le1$ is the grey double cone $x^2=y^2+z^2$, and the model $\rho(c,\phi)=W_{\rm res}+c\,v_\phi v_\phi\T$, $v_\phi=(\cos\phi,\sin\phi)$, which plays LoRA ($c$ for $B$, $v_\phi$ for $A$); its image is the whole cone, so $d=2$. (a)~Zero initialisation, $c_0=0$ and $W_{\rm res}=\theta_0$: the base point $\theta_0$ is the vertex of the cone, and the first-order image $S_1$ (teal) is a single line, one ruling, so the deficit is $1$ and $\theta_0$ is an apex (\thmref{Definition}{I.5}, \thmref{Theorem}{II.2}(c)). (b)~A split (PiSSA-type) initialisation, $c_0\ne0$ and $W_{\rm res}=\theta_0-P$ with $P=c_0v_{\phi_0}v_{\phi_0}\T$: the vertex moves to $\theta_0-P$ (open diamond), $\theta_0$ lies on a smooth sheet of the cone, and $S_1$ is the tangent plane there, so the deficit is $0$. Both values of $\dim S_1$ are Jacobian ranks of the model. (c)~For $m=n=4096$ and $r=16$, each pair of bars compares the expressive dimension $d(M)$ (grey) with $\dim S_1$ (teal); the dashed red box is the gap between them, the first-order deficit. Zero-init LoRA, stratum $\mathrm T_A$, has $\dim S_1=mr=65{,}536$ against $d=r(m+n-r)=130{,}816$; a split initialisation has $\dim S_1=d$ (\thmref{Theorem}{II.2}(b)); HRA's paired initialisation $u_{2i-1}=u_{2i}$, $k=r/2$, has $\dim S_1=kn-\tfrac{k(k+1)}2=32{,}732$ against $d=rn-\tfrac{r(r+1)}2=65{,}400$, so it too starts at an apex (\thmref{Theorem}{II.7}(b)). The script checks each closed form against a Jacobian rank at $(m,n,r)=(11,9,4)$.
